\section{Tangent cones and instanton densities}\label{sec:cones}

The local model \eqref{eq:adhm-germ} is an analytic statement. Identifying
its natural grading with the grading of a metric tangent cone requires
an additional argument. We give that argument before discussing metrics
on the resulting cone.

\subsection{The affine cone and its Reeb action}
For an affine cone, the Reeb action is the circle action associated to
its positive grading. We normalize it so that a holomorphic volume form
in complex dimension $d$ has weight $d$.
On \eqref{eq:adhm-germ}, the standard action multiplies all linear
coordinates, including $B_1,B_2,i,j$, by the same scalar. The holomorphic
symplectic form has weight two, so this action has the stated normalization.

\begin{theorem}\label{thm:cones}
Every point of $\Mn$ has a unique metric tangent cone. Its affine complex
variety is the conical symplectic germ associated to that point.
If $p=(F,\sum_jn_jx_j)$ has stable residual bundle $F$, then this affine
variety, with its normalized Reeb action, is
\begin{equation}\label{eq:cone-algebraic}
 \mathcal C_p=\CC^{\operatorname{ext}^1(F,F)}
                    \times\prod_j\frM_0(r,n_j)
\end{equation}
with the standard action just described.
\end{theorem}
\begin{proof}
The projective symplectic resolution in
Proposition~\ref{prop:algebraic-model} satisfies the hypotheses of
\cite[Theorems~1.1 and 6.1]{NO26}. That result identifies the canonical
conical symplectic germ with the affine variety of the metric tangent
cone. Its polarized smoothing permits the uniqueness theorem of
\cite[Theorem~1.3]{DS17} to be applied. Neither assertion requires the
singularity to be isolated.

It remains to identify the Reeb action at a stable-residual point.
Endow \eqref{eq:cone-algebraic} with the zero-level flat hyperk\"ahler
quotient metric. Its squared radius is a continuous proper homogeneous
function on the affine quotient, by the Kempf--Ness identification.
Relative to an algebraic homogeneous radius its logarithmic ratio is
bounded, since a link is compact. It is a K\"ahler potential on the free
locus and extends plurisubharmonically across the singular set.

Here the free locus is the whole regular locus. One way to check the
necessary simple-root condition is to use the extended ADHM quiver of
dimension vector $(1,n)$. Simple representations exist for $r\ge2$:
take $B_1$ diagonal with distinct eigenvalues, $B_2=0$, let the first two
columns of $i$ be a vector $a$ with all entries nonzero, and let the first
two rows of $j$ be $b,-b$, where all entries of $b$ are nonzero.
Then $ij=0$. No proper coordinate subspace contains $\operatorname{im}i$,
and no nonzero coordinate subspace lies in $\ker j$.
Thus the representation is simple. The simple-root criterion
\cite[Theorem~1.2]{CB01} and the smooth-locus theorem
\cite[Theorem~1.15]{BSc23} identify regular points with simple
representations, equivalently free points modulo the diagonal scalar.
The assertion for the product follows factor by factor.

On this regular locus the quotient metric is hyperk\"ahler and satisfies
the Ricci-flat Monge--Amp\`ere equation. Its potential and volume weight
therefore satisfy the definition of a weak Ricci-flat cone metric in
\cite[Definitions~A.6--A.8]{LWX21}. Corollary~A.4 of that paper makes the
standard cone K-polystable. By \cite[Theorem~3.5]{LX18}, its Reeb valuation
minimizes normalized volume. The minimizer is unique up to scaling \cite[Theorem~2.4]{NO26};
normalizing the volume-form weight removes that scalar. The stable
degeneration theorem, in the form recalled in \cite[Section~2]{NO26},
identifies the first associated graded cone with
\eqref{eq:cone-algebraic}. It is already K-polystable, so the unique
K-polystable degeneration of \cite[Theorem~1.2]{LWX21} makes no further
change. This is the grading of the metric tangent cone.
\end{proof}

\begin{remark}\label{rem:cone-uniqueness}
Theorem~\ref{thm:cones} identifies the affine cone and its Reeb action.
To identify its metric with the flat ADHM quotient metric in full
generality, we would also need the following uniqueness assertion:
\begin{quote}
\textup{(Q)} A Ricci-flat K\"ahler cone metric with locally bounded
potential on a K-polystable cone is unique, for a fixed Reeb action,
up to an automorphism commuting with that action.
\end{quote}
Under \textup{(Q)}, the two cone metrics in question are isometric.
We retain this implication as conditional, in view of the general
uniqueness issue described in \cite[Remark~6.14]{NO26}; compare
\cite[Proposition~4.8]{DS17} and \cite[Theorem~A.12]{LWX21}.
No such additional assertion is needed in
Theorem~\ref{thm:orbifold}, where the tangent metric is explicitly flat
upstairs. It is also unnecessary for the volume densities below.
\end{remark}

\subsection{Densities from the Hilbert character}
For a graded affine variety $C$, write
$H_C(q)=\sum_{m\ge0}\dim\CC[C]_m q^m$ for its Hilbert character.
The normalized metric density at $p$ is
\[
 \vartheta(p)=\lim_{\rho\downarrow0}
 \frac{\Vol(B_\rho(p))}{\omega_{4N}\rho^{4N}}.
\]
The limit exists by volume comparison.

\begin{theorem}\label{thm:density}
At an ideal point $p=(F,\sum_jn_jx_j)$ with stable residual bundle $F$,
\begin{equation}\label{eq:density}
 \vartheta(p)=\prod_j\theta(r,n_j),\qquad
 \theta(r,n)=\lim_{t\downarrow0}t^{2rn}
                    H_{\frM_0(r,n)}(e^{-t}).
\end{equation}
The factor $\theta(r,n)$ is the $n$-instanton coefficient of the
$U(r)$ Nekrasov partition function at
$\varepsilon_1=\varepsilon_2=1$ and $a_1=\cdots=a_r=0$, with
specialization performed after summing and cancelling the rational
localization expression. In particular,
\begin{equation}\label{eq:one-instanton}
 \theta(r,1)=4^{-(r-1)}\binom{2r-2}{r-1}.
\end{equation}
\end{theorem}
\begin{proof}
Noncollapsed volume convergence identifies $\vartheta(p)$ with the vertex
density of the tangent cone. The leading Hilbert-character coefficient of
a normalized Ricci-flat cone is its volume ratio
\cite[Appendix~C]{HeinSun17}. By Theorem~\ref{thm:cones}, the character is
that of \eqref{eq:cone-algebraic}. Characters multiply under products;
the Euclidean factor contributes one. This proves \eqref{eq:density}
without requiring a product decomposition of the tangent metric.

The Gieseker resolution
$\frM(r,n)\to\frM_0(r,n)$ has a rational-singularities target. Its
equivariant Euler characteristic of the structure sheaf is therefore the
character of the affine coordinate ring. The leading term of equivariant
Riemann--Roch is the equivariant integral of one, computed in
\cite[Proposition~4.1 and (4.6)]{NY05}. This is the instanton coefficient
in the statement. Standard linear scaling of the ADHM data agrees on the
quotient with the torus specialization
$\varepsilon_1=\varepsilon_2=1$, $a=0$: the two lifts differ by a central
$\operatorname{GL}(n)$ action. Individual fixed-point terms may have poles
there; their sum, interpreted as this positively graded equivariant
multiplicity, has the asserted finite value.

For $n=1$ the quotient is $\CC^2$ times the minimal nilpotent orbit closure
in $\mathfrak{sl}_r$. The coordinate ring of the minimal orbit closure in degree $2k$ has dimension
\[
 \dim V_{k\vartheta}=\frac{2k+r-1}{r-1}
                         \binom{k+r-2}{r-2}^2,
\]
where $\vartheta$ here denotes the highest root.
The leading coefficient of this polynomial in $k$ is
$2/((r-1)((r-2)!)^2)$ and its degree is $2r-3$.
Since $\sum_{k\ge0}k^m e^{-2kt}\sim m!/(2t)^{m+1}$, the leading
character coefficient is
\[
 \frac{2(2r-3)!}{(r-1)((r-2)!)^2\,2^{2r-2}}
 =4^{-(r-1)}\binom{2r-2}{r-1},
\]
as required.
\end{proof}

For reproducibility, the rational expression used for the higher values
is as follows. Let $\mathbf Y=(Y_1,\ldots,Y_r)$ range over tuples of Young
diagrams with total size $n$. For a box $s=(i,j)$, put
$A_Y(s)=Y_i-j$ and $L_Y(s)=Y'_j-i$, allowing negative values when $s$ is
outside the diagram. Define
\begin{align*}
 E_{\alpha\beta}(\mathbf Y)
 ={}&\prod_{s\in Y_\alpha}
 [a_\beta-a_\alpha-\varepsilon_1L_{Y_\beta}(s)
                     +\varepsilon_2(A_{Y_\alpha}(s)+1)]\\
 &\quad\cdot\prod_{s\in Y_\beta}
 [a_\beta-a_\alpha+\varepsilon_1(L_{Y_\alpha}(s)+1)
                     -\varepsilon_2A_{Y_\beta}(s)].
\end{align*}
Then
\begin{equation}\label{eq:localization}
 Z_n(\varepsilon_1,\varepsilon_2;a)
       =\sum_{|\mathbf Y|=n}\prod_{\alpha,\beta=1}^r
                                  E_{\alpha\beta}(\mathbf Y)^{-1}
\end{equation}
is the tangent-weight formula of \cite[Theorem~2.11]{NY05}.
For example, with $r=2$ and $a_1=a$, $a_2=-a$, writing
$e=\varepsilon_1+\varepsilon_2$ gives
\[
 Z_2=
 \frac{8e^2+\varepsilon_1\varepsilon_2-8a^2}
 {\varepsilon_1^2\varepsilon_2^2(e^2-4a^2)
 ((2\varepsilon_1+\varepsilon_2)^2-4a^2)
 ((\varepsilon_1+2\varepsilon_2)^2-4a^2)}.
\]
Substitution gives $\theta(2,2)=33/324=11/108$.
The same finite sums give Table~\ref{tab:densities}.
The accompanying exact-arithmetic calculation evaluates a generic
one-parameter approach before taking its limit, so it never assigns a
value to an individual singular summand.

\begin{table}[ht]
\centering
\caption{Normalized densities of the affine instanton cones.}\label{tab:densities}
\begin{tabular}{c@{\qquad}ccccc}
\toprule
$(r,n)$ & $(2,1)$ & $(2,2)$ & $(2,3)$ & $(3,1)$ & $(3,2)$\\
\midrule
$\theta(r,n)$ & $1/2$ & $11/108$ & $61/5184$ & $3/8$ & $61/1296$\\
\bottomrule
\end{tabular}
\end{table}

The rank-one value $\theta(1,n)=1/n!$, obtained from
$\Sym^n\CC^2$, is a normalization check outside the rank assumption of the
main theorem. The inequalities
\[
 \frac{11}{108}<\frac14,\qquad
 \frac{61}{5184}<\frac{11}{216},\qquad
 \frac{61}{1296}<\left(\frac38\right)^2
\]
are consistent with lower semicontinuity of metric density when points
collide. Neither the product formula nor this table describes reducible
residual bundles: their full Ext quivers replace
\eqref{eq:cone-algebraic}.
